% =========================================================================
\section{Resumen de resultados del PreLab}\label{sec:resumen}
% =========================================================================
La \cref{tab:resumen} reúne los resultados que deben figurar, enmarcados, en el PreLab manuscrito.

\begin{table}[ht]
\centering\small
\renewcommand{\arraystretch}{1.45}
\caption{Resultados de los tres PreLab con el modelo ideal del AO.}
\label{tab:resumen}
\begin{tabular}{@{}lccc@{}}
\toprule
 & Inversor & No inversor & Inversor generalizado\\
\midrule
Ganancia & $A_V=-\dfrac{R_2}{R_1}$ \ \eqref{eq:inv-Av} & $A_V=1+\dfrac{R_2}{R_1}$ \ \eqref{eq:ni-Av}
         & $H(s)=-\dfrac{Z_2}{Z_1}$ \ \eqref{eq:gen-H}\\
Diseño ($R_1=1\,\mathrm{k}\Omega$) & $R_2=10\,\mathrm{k}\Omega$ \ \eqref{eq:inv-R2} & $R_2=10\,\mathrm{k}\Omega$ \ \eqref{eq:ni-R2}
       & $Z_1=R$, $Z_2=R_F\parallel C$\\
Resultado & $A_V=-10$ & $A_V=+11$ & $H(s)=-\dfrac{R_F}{R}\,\dfrac{1}{1+sR_FC}$ \ \eqref{eq:pb-H}\\
\bottomrule
\end{tabular}
\end{table}

Los tres circuitos se resuelven con las mismas dos reglas del AO ideal: corrientes de entrada nulas
y cortocircuito virtual. La ganancia la fija la red pasiva de realimentación, no el AO. Si la red
contiene un condensador, la misma deducción, hecha con impedancias en el dominio de Laplace, da un
filtro paso bajo: ganancia $R_F/R$ por debajo de $f_c=1/(2\pi R_FC)$ y caída de $-20\dB$/década por encima.
